\section{Conclusion}

This project successfully bridged the theoretical foundations of computer networking with hands-on system implementation. 
By developing a hybrid chat system from the ground up, the project demonstrated a comprehensive understanding of key networking paradigms—Client–Server and Peer-to-Peer (P2P)—and their practical integration into a cohesive, multi-process application. 
The resulting system not only satisfied the technical specifications of the assignment but also offered valuable insights into the challenges of designing reliable distributed systems.

\subsection{Summary of Achievements}

The project achieved all primary objectives outlined in the assignment description, delivering a fully functional and compliant implementation. 
The main accomplishments are summarized below:

\subsubsection{Complete HTTP Server and Session Management}
A multithreaded HTTP server was implemented entirely from scratch, capable of handling request parsing, routing, and static content delivery. 
A robust cookie-based session mechanism was built to maintain user authentication state across the stateless HTTP protocol, meeting the functional requirements of Tasks~1A and~1B.

\subsubsection{Hybrid Client–Server and Peer-to-Peer Architecture}
The system integrates a centralized Client–Server model—used for authentication, peer discovery, and state management—with a decentralized P2P network layer for direct, low-latency communication. 
This architecture fulfills the objectives of Task~2 by demonstrating how hybrid systems can combine reliability with efficiency in real-time data exchange.

\subsubsection{Multi-Process Integration}
All system components—the reverse proxy, the backend application server, and the per-user P2P daemons—were successfully integrated into a unified network environment. 
These independent processes interact through well-defined HTTP and TCP interfaces, illustrating the ability to design, coordinate, and test complex multi-process communication flows, as expected in Task~3.

\subsubsection{Full Compliance with Assignment Constraints}
The entire project adheres strictly to the assignment’s restrictions: no use of external libraries, frameworks, or WebSocket-based communication. 
All networking features were implemented using only the Python standard library, relying on modules such as \texttt{socket}, \texttt{threading}, and \texttt{mimetypes}. 
This compliance required deep engagement with low-level network programming concepts, manual protocol definition, and concurrency management.

\subsection{Lessons Learned}

The development process yielded several critical lessons that extend beyond theoretical understanding and directly contribute to practical engineering competence.

\subsubsection{The Synergy and Challenge of Integrating Protocols}
Building an "HTTP Bridge" that links browser-based HTTP communication with low-level TCP socket exchanges offered invaluable insight into the interplay between networking layers. 
This experience demonstrated how protocol design can be used creatively to overcome real-world constraints—such as browser restrictions—while preserving communication efficiency and correctness.

\subsubsection{The Criticality of Concurrent Programming and State Management}
Implementing a server capable of handling multiple clients concurrently emphasized the challenges of concurrency in network programming. 
Through direct experience with Python’s \texttt{threading} module and synchronization primitives like \texttt{RLock}, the project reinforced best practices for ensuring thread safety and preventing race conditions in shared resources (e.g., the Tracker’s peer registry).

\subsubsection{The Foundational Role of Protocol Design}
Defining custom protocols for P2P connection establishment (\texttt{CONNECT/ACCEPT}) and connection maintenance (\texttt{PING/PONG}) highlighted that clear, well-specified protocols form the backbone of reliable distributed systems. 
A precise protocol definition simplifies implementation, facilitates debugging, and ensures consistent interpretation of events between communicating peers. 
This reinforced the understanding that sound protocol and architectural design are essential first steps in developing scalable and fault-tolerant network applications.

\vspace{0.5em}
In conclusion, the project fully achieved its objectives by successfully transforming core networking principles into a functioning hybrid communication system. 
It demonstrated technical competence in socket programming, concurrency, and protocol design—precisely the outcomes intended by the Computer Networks course. 
Beyond meeting assignment requirements, it provided an invaluable foundation for future exploration of scalable, internet-wide distributed architectures.
